% SEGMENT OLP-0575-S01
% Part: set-theory
% Chapter: ord-arithmetic
% Section: introduction

\documentclass[../../../include/open-logic-section]{subfiles}

\begin{document}
\olfileid{sth}{ord-arithmetic}{intro}

\olsection{அறிமுகம்}

\olref[ordinals][]{chap} இல் வரிசையெண்களின் கோட்பாட்டை உருவாக்கினோம்.
\olref[spine][]{chap} இல், படிநிலைமுறையின் எஞ்சிய பகுதியைக் கட்டமைப்பதற்கான
முதுகெலும்பாக வரிசையெண்களைக் கருதலாம் என்று கண்டோம். ஆனால் அவற்றின் பங்கு
அதோடு முடிவதில்லை. வரிசையெண் எண்கணிதமும் செய்ய வேண்டியுள்ளது.

முன்னரே \olref[ordinals][idea]{sec} இல் $\omega$, $\omega+1$,
$\omega+\omega$ ஆகியவற்றைப் பற்றிப் பேசியபோது இதைச் சுட்டிக்காட்டினோம்.
அப்போது முறைசாராமல் பேசினோம்; இப்போது இதைத் துல்லியமாக விளக்க வேண்டிய
நேரம் வந்துவிட்டது. எனினும், இந்த அத்தியாயத்தில் தத்துவப் பகுதி அதிகம்
இல்லை என்று சொல்ல வேண்டும்: ஒரே காரியத்தை இரண்டு வெவ்வேறு வழிகளில்
செய்யலாம் என்ற (ஓரளவு) சுவாரசியமான கருத்துடன் கூடிய தொழில்நுட்ப
வளர்ச்சிகள்தாம் இங்கு உள்ளன.

\end{document}
